% META: {"id": "TCOMPL-LOG-ME-003", "chapitre": "TCOMPL-LOGARITHME-HISTORIQUE", "type_objet": "methode", "capacites_codes": ["C3"], "methodes": ["M3"], "mode_creation": "ex_nihilo", "sources_inspiration": [], "status": "needs_review"}
% BEGIN-VERIFY
% from sympy import Rational, log, floor
% from math import log as mlog
% q, seuil = Rational(8,10), Rational(1,100)
% n0 = mlog(float(seuil))/mlog(float(q))
% assert 20 < n0 < 21
% assert q**20 > seuil and q**21 < seuil
% END-VERIFY
\begin{fichemethode}{M3}{Determiner un seuil grace au logarithme d'une puissance}

\quandUtiliser{%
  Utiliser cette methode lorsque l'enonce demande le plus petit entier $n$ tel
  qu'une suite geometrique franchisse un seuil : resoudre $q^n < s$ (ou
  $q^n > s$) a l'aide de $\ln(q^n) = n\ln(q)$.
}

\pasApas{%
  \begin{enumerate}
    \item Ecrire l'inequation a resoudre (ex. $u_0\,q^n < s$) et l'isoler sous
      la forme $q^n < s'$ avec $q > 0$ et $s' > 0$.
    \item Appliquer $\ln$ aux deux membres (fonction strictement croissante,
      le sens est conserve) : $\ln(q^n) < \ln(s')$.
    \item Utiliser $\ln(q^n) = n\ln(q)$ pour obtenir $n\ln(q) < \ln(s')$.
    \item Diviser par $\ln(q)$ en surveillant son SIGNE : si $0 < q < 1$,
      $\ln(q) < 0$ et l'inegalite CHANGE de sens :
      $n > \dfrac{\ln(s')}{\ln(q)}$.
    \item Conclure par le plus petit entier qui convient, et le verifier par un
      calcul direct de $q^n$ de part et d'autre.
  \end{enumerate}
}

\exempleRedige{%
  Determiner le plus petit entier $n$ tel que $0{,}8^n < 0{,}01$.

  \medskip
  \commentaireMarge{$\ln$ conserve le sens de l'inegalite.}%
  $0{,}8^n < 0{,}01 \iff n\ln(0{,}8) < \ln(0{,}01)$.

  \medskip
  \commentaireMarge{$\ln(0{,}8) < 0$ : division qui RENVERSE l'inegalite.}%
  Comme $\ln(0{,}8) < 0$ :
  $n > \dfrac{\ln(0{,}01)}{\ln(0{,}8)} \approx 20{,}64$.

  \medskip
  Le plus petit entier qui convient est $n = 21$. Controle :
  $0{,}8^{20} \approx 0{,}0115 > 0{,}01$ et $0{,}8^{21} \approx 0{,}0092 < 0{,}01$.
}

\pieges{%
  \begin{itemize}
    \item \textbf{Piege 1.} Diviser par $\ln(q)$ sans regarder son signe : pour
      $0 < q < 1$, $\ln(q) < 0$ et le sens de l'inegalite se renverse.
    \item \textbf{Piege 2.} Arrondir vers le bas : si $n > 20{,}64$, le premier
      entier est 21, pas 20.
    \item \textbf{Piege 3.} Ecrire $\ln(q^n) = \ln(q)^n$ : la propriete est
      $n\ln(q)$.
  \end{itemize}
}

\verifier{%
  Toujours tester l'entier trouve ET le precedent par calcul direct : $q^{n}$
  doit franchir le seuil, $q^{n-1}$ non.
}

\sEntrainer{\refExos{M3}}

\end{fichemethode}
